Convolutional Neural Networks (CNNs) are primarily driven by local texture features rather than global shape information, unlike human vision. While previous research demonstrated that a ``biomimetic'' training regimen---starting with blurry, achromatic inputs before transitioning to high-resolution, full-color images---can induce a more human-like shape bias, the importance of the specific temporal order of this experience remains unclear. In this study, we investigate the impact of an ``anti-biomimetic'' regimen that reverses this developmental sequence. Using a modified AlexNet architecture trained on a subset of the ImageNet database, we compared standard, biomimetic, and anti-biomimetic protocols. Behavioral analyses reveal that the anti-biomimetic model failed to develop a shape bias and instead exhibited a stronger texture bias than even the standard control network. Furthermore, receptive field characterization showed that reversing the sensory progression prevented the emergence of the magnocellular-like clusters (low spatial frequency, low color sensitivity) observed in biomimetic models. These results suggest the existence of a ``critical period'' for structural reorganization early in training, where the temporal order of sensory experience is the primary driver behind visual pathway organization.